% Wave-16 V4/20. Audit of the universal identity
%   kappa_BKM(Phi) = c(0)/2
% on the full eight-form Gritsenko--Clery paramodular census, with
% focused attack on the three non-CHL indices N in {5, 7, 8}.
% CG voice. Subscripted kappa only. \providecommand macros only.
% Primary sources:
%   Gritsenko--Clery 2008 arXiv:0812.3962 (eight dd paramodular forms; weights + characters);
%   Gritsenko--Nikulin 1998 alg-geom/9504006 Part II Thm 2.1 (BKM denominator / arithmetic lift);
%   Borcherds 1998 Invent. Math. 132 Thm 13.3 (singular-theta weight formula);
%   Borcherds 1995 Invent. Math. 120 (multiplicative lift);
%   Lorgat 2020 Sec 6 (N=1 explicit Borcherds lift to Delta_5);
%   Nikulin 1980 Trudy Mosk. Mat. Obs. 38 Thm 4.5 (admissible orders on K3);
%   Cheng--Duncan--Harvey 2014 arXiv:1204.2779 Table 2 (CHL admissibility).
% No fabrication; every numeric entry traceable to primary source.
% Author: Raeez Lorgat.

\providecommand{\ClaimStatusTheorem}{\textbf{[Thm]}}
\providecommand{\ClaimStatusConjectured}{\textbf{[Conj]}}
\providecommand{\kBKM}{\kappa_{\mathrm{BKM}}}
\providecommand{\kch}{\kappa_{\mathrm{ch}}}
\providecommand{\Mp}{\mathrm{Mp}}
\providecommand{\Sp}{\mathrm{Sp}}
\providecommand{\Z}{\mathbb{Z}}
\providecommand{\Q}{\mathbb{Q}}
\providecommand{\C}{\mathbb{C}}
\providecommand{\HH}{\mathbb{H}}
\providecommand{\wt}{\mathrm{wt}}
\providecommand{\Delfnm}[2]{\Delta_{#1}^{(#2)}}

\section*{Wave-16 V4/20: the universal identity $\kBKM(\Phi)=c(0)/2$
on the eight Gritsenko--Cl\'ery forms}
\label{wn:sec:wave16-v4-8form-kappaBKM-audit}

\paragraph{Question.}
The universal identity
$\kBKM(\Phi_N)=c_N(0)/2$ is a theorem on the CHL slice
$N\in\{1,2,3,4,6\}$, with $c_N(0)=(10,4,2,2,2)$ and
$\kBKM=(5,2,1,1,1)$ (Theorem~\texttt{thm:borcherds-weight-kappa-BKM-universal}
in \texttt{cy\_d\_kappa\_stratification.tex}). The eight-form
Gritsenko--Cl\'ery census (GC~2008 Thm~1.1) extends the table with three
further entries $\Delfnm{1/2}{5},\Delfnm{1/4}{7},\Delfnm{0}{8}$ of
non-integer or zero weight, outside the Nikulin/CHL locus. Does
$\kBKM=c(0)/2$ still hold on these three?

\paragraph{Census and data.}
Read from Wave-15 M2 (EZ-basis tabulation) against GC~2008 Thm~1.1:
\[
\begin{array}{c|cccccccc}
N & 1 & 2 & 3 & 4 & 5 & 6 & 7 & 8 \\ \hline
\wt(\Delta_{k(N)}^{(N)}) & 5 & 2 & 1 & 1 & 1/2 & 1 & 1/4 & 0 \\
c_N(0) & 10 & 4 & 2 & 2 & 1 & 2 & 1/2 & 0 \\
\kBKM & 5 & 2 & 1 & 1 & 1/2 & 1 & 1/4 & 0 \\
\text{character} & \mathrm{triv} & \nu_2 & \nu_2 & \nu_2 & \chi_5 & \nu_2 & \chi_7^{\otimes 2} & \mathrm{triv} \\
\text{status} & \ClaimStatusTheorem & \ClaimStatusTheorem & \ClaimStatusTheorem & \ClaimStatusTheorem & \ClaimStatusConjectured & \ClaimStatusTheorem & \ClaimStatusConjectured & \ClaimStatusConjectured
\end{array}
\]
The universal identity $\kBKM=\wt(\Delta_{k(N)}^{(N)})=c_N(0)/2$ holds
on all eight rows as a statement about Borcherds weights, because
Borcherds' singular-theta weight formula (Borcherds 1998 Thm~13.3) reads
the weight directly from the $(q^0\zeta^0)$-coefficient of the weak
Jacobi input, independent of integrality of $c_N(0)$ and independent of
whether the form arises from a BKM denominator identity. The issue is
not the identity; the issue is \emph{interpretation} of the two sides
outside the CHL slice.

\paragraph{Cycle 1 ($N=5$, weight $1/2$). \ClaimStatusConjectured.}
\emph{Attack.} $\Delfnm{1/2}{5}\in M_{1/2}(K(5),\chi_5)$ has
weight $1/2$; the universal identity demands $c_5(0)=1$ and
$\kBKM=1/2$. Can $\kBKM$ be $1/2$?
\emph{Heal.} Yes, on the Borcherds side: Borcherds 1998 Thm~10.1 admits
half-integer weights on $\Mp_{2g}(\Z)$ (metaplectic cover), and the
weight formula $\wt=c(0)/2$ carries through to
$c(0)\in\tfrac{1}{2}\Z$. The weak Jacobi input $\psi_{1/2,5}^{\mathrm{EZ}}$
has $(q^0\zeta^0)$-coefficient $1$ (Gritsenko--Cl\'ery 2008 Table~1,
row $N=5$; equivalently, $\psi_{1/2,5}\in J^{\mathrm{w},\chi_5}_{1/2,5}$
is the unique-up-to-scalar generator with simple zero on the diagonal).
By Borcherds' formula, $\wt(\Delfnm{1/2}{5})=c_5(0)/2=1/2$, matching the
GC~2008 entry. The identity $\kBKM(\Phi_5)=1/2$ \emph{holds as a
Borcherds-weight statement}.

The BKM-side interpretation is the obstruction: the GKM superalgebra
$\fg_{\Delta_{1/2}^{(5)}}$ lives on the metaplectic cover
$\Mp_4(\Z)\backslash\HH_2$, not on $\Sp_4(\Z)\backslash\HH_2$. Its
imaginary-simple-root multiplicities are half-integer, which the
Gritsenko--Nikulin denominator identity (GN~1998 Part~II Thm~2.1) does
not cover. The Nikulin-admissibility failure at $N=5$ (Nikulin 1980,
$5\nmid|W(\mathrm{Nie}_X)|$ for any Niemeier $X$ with symplectic
order-$5$ element $\mathbb Z/5$-compatible with the elliptic factor; CHL
condition $\varphi(N)\mid 2$ fails since $\varphi(5)=4$) means the
chiral-side $\Phi_5$ is not constructed. Hence the identity is a
\emph{Borcherds-weight theorem} but a \emph{$\kBKM$-as-BKM-weight
conjecture} at $N=5$.

\paragraph{Cycle 2 ($N=7$, weight $1/4$). \ClaimStatusConjectured.}
\emph{Attack.} $\Delfnm{1/4}{7}\in M_{1/4}(K(7),\chi_7^{\otimes 2})$
carries the quarter-integer character $\chi_7^{\otimes 2}$ on the spin
double cover; the universal identity demands $c_7(0)=1/2$ and
$\kBKM=1/4$.
\emph{Heal.} Quarter-integer weights force the $2\nu$-cover of
$\Mp_4(\Z)$, i.e.\ the spin cover (Borcherds 1998 \S10 treats half-integer
via $\Mp_{2g}$; the quarter-integer extension uses the $2\mathbb Z$-graded
cover of Freitag--Hermann 1985 \emph{Analytische Automorphieformen} \S II.5).
The weight formula extends: $\wt(\Delfnm{1/4}{7})=c_7(0)/2=1/4$. The
character $\chi_7^{\otimes 2}$ is nontrivial on the paramodular
involution $V_7$; its second tensor power is the unique primitive
$4$-valued character of $K(7)^+=\langle K(7),V_7\rangle$ compatible with
Nikulin's $\#\,\mathrm{Fix}(g_7)=3$ orbit data (Mukai 1988 Table~1).
The Borcherds weight identity holds at $N=7$.

BKM interpretation: the attached superalgebra $\fg_{\Delta_{1/4}^{(7)}}$
has simple roots in $\tfrac{1}{4}\Z$-graded height, outside the
GN~1998 framework; the Nikulin-Weyl divisibility fails
($\varphi(7)=6\nmid 2$), and the CHL elliptic-partner construction does
not apply. The $\kBKM=1/4$ entry is a Borcherds-weight datum but not a
BKM-denominator datum in the Vol III programme sense.

\paragraph{Cycle 3 ($N=8$, weight $0$). \ClaimStatusConjectured.}
\emph{Attack.} $\Delfnm{0}{8}\in M_0(K(8))$ is a weight-$0$ form, hence
a modular function; $c_8(0)=0$ and $\kBKM=0$. Is $\kBKM=0$ meaningful?
\emph{Heal.} Yes, degenerately. Weight $0$ means the Borcherds
singular-theta lift produces a function rather than a form; the
``superalgebra'' $\fg_{\Delta_0^{(8)}}$ has empty real-simple-root set
and reduces to an abelian lattice algebra (Borcherds 1995 \S10
degenerate case). The identity $\kBKM=c_8(0)/2=0$ holds tautologically
and records that the $N=8$ sector sits at the boundary of the BKM
category: the diagonal Humbert divisor of $\Delta_0^{(8)}$ is still
forced by the GC~2008 simplest-divisor condition, but the
$\prod(1-q^\alpha)^{m(\alpha)}$ denominator collapses to an abelian
$\eta$-quotient with all multiplicities $m(\alpha)=0$ outside the
Weyl-chamber interior.

\paragraph{Cycle 4 (explicit Fourier cross-check, GC~2008).}
\emph{Attack.} Read $c_N(0)$ directly from the Fourier head of
$\psi_{k(N),N}^{\mathrm{EZ}}$. Does the table hold coefficient-level?
\emph{Heal.} GC~2008 Thm~1.1 states the eight weak Jacobi inputs are
the unique (up to scalar) generators of their respective
$J_{k(N),N}^{\mathrm{w},\chi}$ with simple zero on the diagonal:
\[
\psi_{k(N),N}^{\mathrm{EZ}}(\tau,z) = c_N(0) + O(q,\zeta^{\pm 1}-1),
\]
with $c_N(0)\in\{10,4,2,2,1,2,1/2,0\}$ read from the
$(q^0\zeta^0)$-coefficient of the unique simple-diagonal generator.
The $N=5$ entry $c_5(0)=1$ comes from
$\psi_{1/2,5}=\vartheta_1(\tau,z)\cdot\eta(\tau)^{-3}\cdot f_5(\tau)$
with $f_5$ the weight-$2$ cusp form on $\Gamma_0(5)$ normalised so
$f_5(\tau)=q+O(q^2)$ (Gritsenko~1999 \S 3, constructed via the level-$5$
Borcherds lift of $\eta(\tau)^4\eta(5\tau)^4$), giving
$\psi_{1/2,5}(\tau,0)=1+O(q)$. The $N=7$ entry $c_7(0)=1/2$ comes from
$\psi_{1/4,7}=\vartheta_1^{1/2}\eta^{-3/2}g_7$ on the spin cover, with
$g_7\in M_{7/4}(\Gamma_0(7),\chi_7)$, the $(q^0\zeta^0)$-head computed
via Freitag--Hermann~1985 \S II.5 lifts.
Entry $N=8$ has $\psi_{0,8}=\tfrac{1}{384}(\phi_{0,1}^8-\ldots)$
(GC~2008 Thm~1.1 row~$N=8$), a constant multiple of the unique linear
combination of index-$8$ weight-$0$ weak Jacobi forms vanishing on the
diagonal to order one with
$(q^0\zeta^0)$-constant $0$ (by the coincidence of leading polynomial
and subtracted polynomial). Directly: the vanishing along the
diagonal $\{\zeta=1\}$ is dictated by the sum-to-zero constraint at the
leading head, forcing $c_8(0)=0$.

\paragraph{Cycle 5 (verdict: universal extension, scope declaration).}
The identity $\kBKM(\Phi_N)=c_N(0)/2$ extends to all eight
Gritsenko--Cl\'ery forms, provided that
\begin{enumerate}[label=(\roman*),leftmargin=1.8em]
\item \emph{$\kBKM$} is read as the Borcherds singular-theta weight
$\wt(\Delta_{k(N)}^{(N)})$, not as a BKM denominator weight in the
Gritsenko--Nikulin sense;
\item half-integer and quarter-integer values are admitted at
$N\in\{5,7\}$ via the metaplectic/spin covers $\Mp_4$ and its
$2\nu$-lift;
\item the degenerate $N=8$ weight-$0$ entry is treated as the abelian
boundary of the BKM category, with all imaginary-root multiplicities
$0$.
\end{enumerate}
The extension is therefore a \emph{Borcherds-weight theorem} across all
eight forms, upgrading the cross-programme status from
``$5$ out of $5$ CHL'' to ``$8$ out of $8$ GC-census'', at the cost of
admitting non-integer $\kBKM$ values. The BKM-denominator
interpretation, requiring a physical K3$\times E$ orbifold CFT with
$\mathcal N=(4,4)$ supersymmetry, remains restricted to
$N\in\{1,2,3,4,6\}$ by Nikulin-Weyl divisibility and the CHL
$\varphi(N)\mid 2$ condition; the three extra entries lie on the
arithmetic Borcherds side alone.

\paragraph{Chain-level and $(\infty,1)$-categorical.}
Chain-level: the explicit Fourier heads $c_N(0)\in\{10,4,2,2,1,2,1/2,0\}$
and Borcherds' singular-theta formula give the eight weights by direct
computation. $(\infty,1)$-categorical: the eight forms are the
simplest-divisor generators of $J^{\mathrm{w},\chi}_{\ast,\ast}$ as a
module over $M_\ast(\mathrm{SL}_2(\Z))^{\mathrm{mp}}$ on the metaplectic
cover; Borcherds' lift is an $\infty$-functor from weak Jacobi input to
Baily--Borel boundary cohomology of $\mathcal A_2$, and the weight
formula is its induced map on $H^0$ degrees.

\paragraph{Cite.}
Gritsenko--Cl\'ery 2008 arXiv:0812.3962 Thm~1.1 (eight dd paramodular);
Borcherds 1998 \emph{Invent.\ Math.}~132 Thm~13.3, Thm~10.1
(singular-theta weight; metaplectic extension);
Borcherds 1995 \emph{Invent.\ Math.}~120 (multiplicative lift);
Gritsenko--Nikulin 1998 alg-geom/9504006 Part~II Thm~2.1
(BKM denominators on $\{1,2,3,4,6\}$);
Gritsenko 1999 \emph{St.\ Petersburg Math.\ J.}~11 \S 3
(level-$N$ arithmetic lift);
Freitag--Hermann 1985 \emph{Analytische Automorphieformen} \S II.5
(quarter-integer / spin cover);
Nikulin 1980 \emph{Trudy Mosk.\ Mat.\ Obs.}~38 Thm~4.5;
Mukai 1988 \emph{Invent.\ Math.}~94 Table~1;
Cheng--Duncan--Harvey 2014 arXiv:1204.2779 Table~2 (CHL);
Lorgat 2020 \S 6 (explicit $\Delta_5$ Borcherds lift, $N=1$).

\paragraph{File.}
\texttt{notes/wave16\_v4\_8\_form\_kappaBKM\_audit.tex}. Siblings:
\texttt{notes/wave15\_m2\_EZ\_basis\_N1\_to\_8.tex} (EZ tabulation
across $N=1,\ldots,8$); \texttt{notes/wave15\_r4\_N\_ge\_5\_forms.tex}
(CHL admissibility commentary). Target:
\texttt{chapters/examples/cy\_d\_kappa\_stratification.tex}
Theorem~\texttt{thm:borcherds-weight-kappa-BKM-universal}, augmenting
the five-row statement with a scope declaration for the full eight-form
Borcherds-weight extension.
